\subsection{Methodological Limitations}

\textbf{Algorithm dependency}: Results assume edge-weighted recursive bisection. Future algorithms may improve success rates, potentially lowering threshold. However, fundamental proportionality constraints (cannot create $X$\% MM districts from $<X$\% minority base) ensure threshold persists.

\textbf{Single optimization run}: Each configuration tested once. METIS introduces randomness; multiple runs might yield slightly different results. However, 645-configuration comprehensive study (43 states × 15 configs) provides robust statistical patterns validated across diverse contexts.

\textbf{Binary success metric}: We use 50\% as MM threshold. Some legal contexts accept 45--50\% as sufficient. Lowering threshold to 45\% would reduce state-level requirement to approximately 40\% (not 42\%).

\subsection{Geographic Limitations}

\textbf{Clustering metrics}: Moran's I captures global spatial autocorrelation but not local complexities. Other metrics (Geary's C, LISA) might provide additional insights.

\textbf{Contiguity requirement}: Analysis assumes districts must be geographically contiguous. Relaxing this constraint (allowing non-contiguous districts) would lower threshold, but violates standard redistricting principles.

\textbf{Compactness vs VRA tradeoff}: Our edge-weighted approach balances both. Abandoning compactness entirely might enable VRA compliance at lower state percentages, but creates gerrymandered districts.

\subsection{Data Limitations}

\textbf{Census resolution}: Analysis uses census tracts as atomic units. Finer resolution (blocks) might enable better minority concentration, marginally lowering threshold. However, tracts provide computationally feasible resolution for METIS optimization across 43 states.

\textbf{Temporal stability}: Analysis uses 2020 Census demographics. Threshold patterns likely stable across decades absent major demographic shifts, but periodic revalidation as new census data becomes available would strengthen findings.

\subsection{Policy Limitations}

\textbf{Coalition districts}: Analysis considers single minority group (non-white vs white). Multi-racial coalition districts might achieve VRA goals at lower single-group percentages.

\textbf{Political context ignored}: We assess purely demographic/geographic feasibility. Political factors (partisan composition, incumbency) affect real-world redistricting.

\textbf{Proportionality assumption}: We target MM proportion matching state minority percentage. Courts might require different proportions case-specifically.

\subsection{Future Research Directions}

\textbf{Temporal analysis}: Compare thresholds across 2000, 2010, 2020 census years to assess longitudinal stability and demographic trend impacts.

\textbf{Sub-group analysis}: Examine thresholds for Hispanic, Asian, Native American communities separately.

\textbf{Alternative algorithms}: Test whether emerging optimization methods (deep learning, genetic algorithms) lower threshold significantly.

\textbf{Coalition analysis}: Study feasibility of multi-racial coalition districts in states below single-group threshold.

\textbf{Synthetic geography}: Generate controlled synthetic states with varying minority percentages and clustering to isolate causal factors.
